\subsection{Using Special Forms} 
We now have five special forms of polynomials, that we can use to help us factor.
\begin{fact}[Perfect Square Trinomials]
\centering\(
(A+B)^2=A^2+2AB+B^2
\qquad
(A-B)^2=A^2-2AB+B^2
\)
\end{fact}
\begin{fact}[Difference of Squares]
\centering\(
A^2-B^2=(A-B)(A+B)
\)
\end{fact}
\begin{fact}[Sum and Difference of Cubes]
\centering\(
A^3-B^3=(A-B)(A^2+AB+B^2)
\qquad
A^3+B^3=(A+B)(A^2-AB+B^2)
\)
\end{fact}
How do we use this? We should be in the habit, when factoring, of being on the lookout for these special forms.
\begin{itemize}
\item When factoring a trinomial, two \makeblank terms should ``jump out'' at us, and we should check for a perfect square.
\item When factoring a \makeblank, we can \emph{only} factor if we have one of the special forms!
\end{itemize}
\begin{Example}
For each binomial, determine whether it is a special form, and if so, factor it.
\begin{lpic}{}{13}
\foreach \y/\exp in {1/x^4-1,2/8y^6+343,3/5x^4-1,4/25y^2+16,5/x^3-25,6/x^6-64}
	\regtextleft{0}{-\y*2+1}{$\exp$};
\end{lpic}
\end{Example}
\begin{Note}
If a binomial is \emph{both} a difference of squares and of cubes, we'll factor it first as a difference of \makeblank[1.5in].
\end{Note}
